% =============================================================================
% Chapter template (rename this file to e.g. 02_first_paper.tex)
% -----------------------------------------------------------------------------
% This file shows the structure of a typical paper-chapter in a Monash
% thesis-by-publication. The conventional structure is:
%
%   1. \chapter{...} title
%   2. \chaptermark{...} (short title for running header)
%   3. \label{...} for cross-referencing
%
%   4. \section*{Linking text} — a short bridging essay (1,500–3,000 words)
%      that situates this chapter in the thesis's overall argument and
%      connects it to the chapter before and after. Linking texts are
%      examiner-facing, written in the thesis's first-person voice.
%      Not part of the published paper.
%
%   5. \section*{Abstract} — copied verbatim from the published paper
%   6. \section{Introduction} — the published paper begins here
%   7. ... rest of the paper as published
%
% NOTE: For unpublished or in-progress chapters, omit the abstract and
% structure the chapter as a standalone piece.
% =============================================================================

\chapter{Your chapter title}
\chaptermark{Short title for running header}
\label{ch:chapter-template}

\section*{Linking text}

% The linking text bridges from the previous chapter to this one. A typical
% structure:
%   - Brief summary of what the previous chapter established
%   - Statement of the new question this chapter takes up
%   - Explanation of why this is the next move in the thesis's arc
%   - Brief preview of what this chapter will do
%
% Aim for ~2,000 words. The linking texts are where examiners will look
% to assess whether the thesis hangs together as a coherent whole rather
% than a collection of papers.

[The chapter that just concluded ... the natural next question is ... this
chapter takes up ... and proceeds as follows ...]

\clearpage

\section*{Abstract}

% If this chapter is a published paper, paste the abstract verbatim.
% Add a keywords line below if the published version had one.

[Abstract of the published paper goes here.]

\textbf{Keywords:} keyword1; keyword2; keyword3.

\clearpage

\section{Introduction}
\label{sec:chapter-template-intro}

[The published paper begins here.]

\section{Methods}
\label{sec:chapter-template-methods}

[Methods section.]

\section{Results}
\label{sec:chapter-template-results}

[Results section. Reference figures with \Cref{fig:example-figure} and
tables with \Cref{tab:example-table}.]

% Figure example. Place image files in figures/ and reference by filename.
% Use [htbp] (here, top, bottom, page) for flexible placement.
% \begin{figure}[htbp]
%     \centering
%     \includegraphics[width=0.7\linewidth]{example_figure.png}
%     \caption{Caption goes here. Make captions self-explanatory; the
%     reader should be able to understand the figure without referring
%     back to the text.}
%     \label{fig:example-figure}
% \end{figure}

% Table example.
% \begin{table}[htbp]
%     \centering
%     \begin{tabular}{lrr}
%         \toprule
%         Condition & Mean & SD \\
%         \midrule
%         A & 1.23 & 0.45 \\
%         B & 6.78 & 0.91 \\
%         \bottomrule
%     \end{tabular}
%     \caption{Table caption.}
%     \label{tab:example-table}
% \end{table}

\section{Discussion}
\label{sec:chapter-template-discussion}

[Discussion section.]
